\documentclass[../../main.tex]{subfiles}% 注意这里的文件路径不能用 ./main.tex ,否则用latexmk编译子文件会报错
\graphicspath{{\subfix{./image/}}} % 指定图片目录,后续可以直接使用图片文件名
% 注意这里的文件路径不能用 ../../image/ ,否则用latexmk编译子文件会报错

% 例如:
% \begin{figure}[H]
% \centering
% \includegraphics[scale=0.3]{图.png}
% \caption{}
% \label{figure:图}
% \end{figure}
% 注意:上述\label{}一定要放在\caption{}之后,否则引用图片序号会只会显示??.

\begin{document}

\section{方程的根式解}

\begin{definition}
设$K$是域$F$的扩域，若存在中间域序列
\begin{align}
F\subset F_1\subset\cdots\subset F_i\subset F_{i+1}\subset\cdots\subset F_m=K,\label{eq:radical_ext_seq}
\end{align}
使得$F_{i+1}$是$x^{n_{i+1}}-a_{i+1}\in F_i[x]$的分裂域，则称$K$是$F$的\textbf{根式扩张}；当$m=1$时，$K$称为$F$的\textbf{单根式扩张}。
\end{definition}

\begin{definition}
设$F$是域，$f(x)\in F[x]$。若存在$F$的根式扩张$K$包含$f(x)$的分裂域，则称方程$f(x)=0$（对$F$）\textbf{可用根式解}。
\end{definition}

\begin{theorem}
设$K$是域$F$的根式扩张，则存在$F$的正规根式扩张$\overline{K}\supseteq K$。
\end{theorem}
\begin{proof}

对根式扩张的链长$m$用数学归纳法。
当$m=1$时，$K$为单根式扩张，即$K$是$x^{n_1}-a_1\in F[x]$的分裂域，故为$F$的正规扩张。

假设结论对链长$m-1$成立。$F_{m-1}$是$F$的根式扩张，由归纳假设，存在$F$的正规根式扩张$\overline{E}\supseteq F_{m-1}$。
$K=F_m$是$F_{m-1}$的单根式扩张，故$K$是$x^n-\beta\in F_{m-1}[x]$的分裂域。设$\zeta$为$n$次单位根，$\theta$是$x^n-\beta$的根，则$K=F_{m-1}(\zeta,\theta)$。
记$\mathrm{Irr}(\beta,F)=\prod\limits_{i=1}^r(x-\beta_i)$，则$\beta=\beta_1\in\overline{E}$，故$\beta_i\in\overline{E}\ (1\leqslant i\leqslant r)$。设$\alpha_i$是$x^n-\beta_i$的根，取$\alpha_1=\theta$，令
\begin{align*}
\overline{K}=\overline{E}(\zeta,\alpha_1,\dots,\alpha_r)\supseteq F_{m-1}(\zeta,\theta)=K,
\end{align*}
存在中间域链
\begin{align*}
\overline{E}=\overline{E}_0\subseteq\overline{E}_1\subseteq\overline{E}_2\subseteq\cdots\subseteq\overline{E}_r=\overline{K},
\end{align*}
其中$\overline{E}_i=\overline{E}_{i-1}(\zeta,\alpha_i)$，故$\overline{E}_i$是$x^n-\beta_i\in\overline{E}_{i-1}[x]$的分裂域，$\overline{K}$是$\overline{E}$的根式扩张，也是$F$的根式扩张。

下证$\overline{K}$是$F$的正规扩张。$\zeta,\alpha_i\in\overline{K}$，故$x^n-\beta_i$的根均在$\overline{K}$中。记
\begin{align*}
g(x)=\prod\limits_{i=1}^r(x^n-\beta_i)\in F[x]\subseteq\overline{E}[x],
\end{align*}
则$\overline{K}$是$g(x)\in\overline{E}[x]$的分裂域。又$\overline{E}$是$F$的正规扩张，故$\overline{E}$为某$f(x)\in F[x]$的分裂域，因此$\overline{K}$是$f(x)g(x)\in F[x]$的分裂域，即$\overline{K}$是$F$的正规根式扩张且$\overline{K}\supseteq K$.

\end{proof}

\begin{theorem}
若域$F$上的多项式$f(x)$满足方程$f(x)=0$可用根式解，则其伽罗瓦群$G(f(x),F)$为可解群。
\end{theorem}
\begin{proof}

设$E$为$f(x)$的分裂域，由$f(x)=0$可用根式解，存在$F$的根式扩张$K\supseteq E$。由前一定理，不妨设$K$为$F$的正规扩张，且根式扩张链为\eqref{eq:radical_ext_seq}。对链长$m$用数学归纳法。

当$m=1$时，$K$是$x^{n_1}-a_1\in F[x]$的分裂域。设$\zeta_1$为$n_1$次单位根，$\theta_1$为$x^{n_1}-a_1$的根，则
\begin{align*}
F\subseteq F(\zeta_1)\subseteq F(\zeta_1,\theta_1)=K.
\end{align*}
$\mathrm{Gal}(F(\zeta_1)/F)$为Abel群，故可解；$\mathrm{Gal}(F(\zeta_1,\theta_1)/F(\zeta_1))$为循环群。由伽罗瓦基本定理，
\begin{align*}
\mathrm{Gal}(K/F)/\mathrm{Gal}(K/F(\zeta_1))\cong\mathrm{Gal}(F(\zeta_1)/F),
\end{align*}
故$\mathrm{Gal}(K/F)$可解。

假设结论对链长$m-1$成立，则$\mathrm{Gal}(K/F_1),\mathrm{Gal}(F_1/F)$均可解。$F_1$是$F$的正规扩张，故
\begin{align*}
\mathrm{Gal}(K/F)/\mathrm{Gal}(K/F_1)\cong\mathrm{Gal}(F_1/F),
\end{align*}
因此$\mathrm{Gal}(K/F)$可解。

$E,K$均为$F$的正规扩张，故
\begin{align*}
\mathrm{Gal}(K/F)/\mathrm{Gal}(K/E)\cong\mathrm{Gal}(E/F),
\end{align*}
由可解群的商群、子群可解，得$\mathrm{Gal}(E/F)=G(f(x),F)$可解.

\end{proof}

\begin{corollary}
设$x_1,x_2,\dots,x_n$为域$F$上的不定元，$p_1,p_2,\dots,p_n$为其初等对称多项式。当$n\geqslant5$时，$g(x)=\prod\limits_{i=1}^n(x-x_i)$在$F(p_1,p_2,\dots,p_n)$上不可用根式解。
\end{corollary}
\begin{proof}

由已知结论，$G(g(x),F(p_1,\dots,p_n))\cong S_n$。$n\geqslant5$时对称群$S_n$非可解群，故$g(x)=0$不可用根式解.

\end{proof}

\begin{corollary}
若$f(x)\in\mathbb{Q}[x]$，$\deg f(x)\geqslant5$，则方程$f(x)=0$在$\mathbb{Q}$上不一定能用根式解。
\end{corollary}
\begin{proof}

对奇素数$p>3$，存在$h(x)\in\mathbb{Q}[x]$使得$G(h(x),\mathbb{Q})\cong S_p$。$p\geqslant5$时$S_p$非可解群，故$h(x)=0$在$\mathbb{Q}$上不可根式解，因此次数不低于5的有理系数多项式方程不一定可用根式解.

\end{proof}

\begin{theorem}
若域$F$上多项式$f(x)$的伽罗瓦群$G(f(x),F)$为可解群，则$f(x)=0$可用根式解。
\end{theorem}
\begin{proof}

记$G_0=G(f(x),F)$，$E$为$f(x)$的分裂域。$G_0$可解，故存在次正规序列
\begin{align*}
G_0\supset G_1\supset\cdots\supset G_s=\{\mathrm{id}\},
\end{align*}
满足$G_i/G_{i+1}$为循环群。由伽罗瓦基本定理，存在中间域序列
\begin{align*}
F_0=F\subset F_1\subset\cdots\subset F_s=E,\quad F_i=\mathrm{Inv}\,G_i,\ 1\leqslant i\leqslant s.
\end{align*}
对$s$用数学归纳法，证明存在$F$的根式扩张$K\supseteq E$。

当$s=1$时，$\mathrm{Gal}(E/F)=G_0$为循环群，设阶为$n$，$\theta$为$n$次单位根。令$K=E(\theta)$，则$F\subseteq F(\theta)\subseteq K$，$\mathrm{Gal}(K/F(\theta))$为循环群，故$K$是$F(\theta)$的单根式扩张，$F(\theta)$是$F$的单根式扩张，即$K$是$F$的根式扩张且$K\supseteq E$。

假设结论对$s-1$成立。$G_1\triangleleft G_0$，故$F_1$是$F$的正规扩张，且
\begin{align*}
\mathrm{Gal}(F_1/F)\cong G_0/G_1
\end{align*}
为循环群，故存在$F$的正规根式扩张$\overline{F}_1\supseteq F_1$。
设$\overline{E}$为$f(x)$与$\overline{F}_1$定义多项式的分裂域，则$\overline{F}_1\supseteq E$。记$H=\mathrm{Gal}(\overline{E}/\overline{F}_1)\subseteq G_1$，存在$H$的次正规序列
\begin{align*}
H=H_0\supset H_1\supset\cdots\supset H_{s-1}=\{\mathrm{id}\},\quad H_i=G_{i+1}\cap H,\ 1\leqslant i\leqslant s-1,
\end{align*}
满足$H_i/H_{i+1}\cong H_iG_{i+1}/G_{i+1}$为循环群。由归纳假设，存在$\overline{F}_1$的正规根式扩张$K\supseteq\overline{E}\supseteq E$，故$K$也是$F$的根式扩张.

\end{proof}

\begin{corollary}
若$f(x)\in F[x]$，$\deg f(x)\leqslant4$，则$f(x)=0$可用根式解。
\end{corollary}
\begin{proof}

设$E$为$f(x)$的分裂域，则$\mathrm{Gal}(E/F)$同构于$S_4$或其子群。$S_4$及其子群均为可解群，故$G(f(x),F)$可解，因此$f(x)=0$可用根式解.

\end{proof}

\begin{remark}
一元多项式方程能否用根式解，等价于其伽罗瓦群是否为可解群；次数不超过4的多项式方程一定可用根式解，次数不低于5的多项式方程不一定可用根式解，且存在不可根式解的高次方程。
\end{remark}

\begin{example}
设$F$为域，不定元$a_1,a_2$的初等对称多项式为$\theta_1=a_1+a_2,\ \theta_2=a_1a_2$。在$F(\theta)[x]$中，$f(x)=(x-a_1)(x-a_2)=x^2-\theta_1x+\theta_2$，其分裂域为$F(a_1,a_2)$，且$\mathrm{Gal}(F(a_1,a_2)/F(\theta_1,\theta_2))=S_2$。
令$D=a_1-a_2$，则$D\notin F(\theta_1,\theta_2)$，且$D^2\in F(\theta_1,\theta_2)$，$\mathrm{Irr}(D,F(\theta_1,\theta_2))=x^2-D^2=x^2-(\theta_1^2-4\theta_2)$。取$D=\sqrt{\theta_1^2-4\theta_2}$，结合$a_1+a_2=\theta_1,\ a_1-a_2=D$，得
\begin{align*}
a_1=\frac{\theta_1+\sqrt{\theta_1^2-4\theta_2}}{2},\quad a_2=\frac{\theta_1-\sqrt{\theta_1^2-4\theta_2}}{2}.
\end{align*}

特别地，$x^2+x+1$的根为三次单位原根：
\begin{align*}
\zeta=\zeta_1=\frac{1}{2}(-1+\sqrt{-3}),\quad \zeta^2=\overline{\zeta}=\zeta_2=\frac{1}{2}(-1-\sqrt{-3}).
\end{align*}
\end{example}

\begin{example}
设域$E$包含三次单位原根，不定元$a_1,a_2,a_3$的初等对称多项式为
\begin{align*}
\theta_1=a_1+a_2+a_3,\quad \theta_2=a_1a_2+a_1a_3+a_2a_3,\quad \theta_3=a_1a_2a_3.
\end{align*}
令$F=E(\theta_1,\theta_2,\theta_3)$，则三次多项式
\begin{align*}
(x-a_1)(x-a_2)(x-a_3)=x^3-\theta_1x^2+\theta_2x-\theta_3\in F[x],
\end{align*}
其分裂域为$K=F(a_1,a_2,a_3)=E(a_1,a_2,a_3)\supseteq F$，且$\mathrm{Gal}(K/F)=S_3$。
记$x_i=a_i-\frac{\theta_1}{3}\ (i=1,2,3)$，则
\begin{align*}
g(x)=\left(x-a_1+\frac{\theta_1}{3}\right)\left(x-a_2+\frac{\theta_1}{3}\right)\left(x-a_3+\frac{\theta_1}{3}\right)=x^3+px+q,
\end{align*}
其中$p=\theta_2-\frac{\theta_1^2}{3},\ q=-\frac{2}{27}\theta_1^3+\frac{1}{3}\theta_1\theta_2-\theta_3$，故$G(g(x),F)=\mathrm{Gal}(K/F)=S_3$。

$S_3$的正规序列为$S_3\supset A_3\supset\{\mathrm{id}\}$，对应不变子域满足
\begin{align*}
F\subseteq F(\Delta)\subset K,
\end{align*}
其中$\Delta=(x_1-x_2)(x_2-x_3)(x_1-x_3)$，满足$\Delta^2=D=-4p^3-27q^2$，即$F(\Delta)=F(\sqrt{D})$。
此时$\mathrm{Gal}(K/F(\Delta))=A_3=\langle\sigma\rangle$，其中$\sigma(x_1)=x_2,\ \sigma(x_2)=x_3,\ \sigma(x_3)=x_1$。

取$\xi\in K$满足$\xi^3\in F(\sqrt{D})$且$K=F(\sqrt{D})(\xi)$，可取
\begin{align*}
\xi=x_1+x_2\zeta+x_3\overline{\zeta},\quad \xi^3=-\frac{27}{2}q+\frac{3\sqrt{-3}}{2}\Delta,
\end{align*}
或
\begin{align*}
\xi'=x_1+x_2\overline{\zeta}+x_3\zeta,\quad \xi'^3=-\frac{27}{2}q-\frac{3\sqrt{-3}}{2}\Delta.
\end{align*}
联立方程组
\begin{align*}
\begin{cases}
x_1+x_2+x_3=0,\\
x_1+\zeta x_2+\overline{\zeta}x_3=\xi,\\
x_1+\overline{\zeta}x_2+\zeta x_3=\xi',
\end{cases}
\end{align*}
解得
\begin{align*}
x_1=\left(-\frac{q}{2}+\sqrt{\frac{q^2}{4}+\frac{p^3}{27}}\right)^{\frac{1}{3}}+\left(-\frac{q}{2}-\sqrt{\frac{q^2}{4}+\frac{p^3}{27}}\right)^{\frac{1}{3}},
\end{align*}
同理可得$x_2,x_3$。
\end{example}



\end{document}